% French PDF pp. 49--54; proof of 4.6.1 continues from en-08.
Taking \(Y\) representable first, one deduces an isomorphism
\(Q\simeq F\). Taking arbitrary \(Y\) next, one deduces the
announced form of \(\Hom(Y,F)\). Finally considering the
canonical morphism \(X\to Q\), one sees at once that it
corresponds to the sub-\(X\)-sheaf \(R\) of \(X\times X\),
which completes the proof.
\end{proof}

\begin{corollary}{4.6.2}\label{IV.4.6.2}
Let \(G\) be any subfunctor of \(F\) such that
\(\Hom(X,G)\subset\Hom(X,F)\) contains \(R\).
Then the canonical morphism \(p:X\to F\) factors through \(G\).
Since \(p\) is covering (4.4.9 and 4.4.3), it follows that
\(G\) is a refinement of \(F\).
In particular, every subsheaf \(G\) of \(F\) satisfying the
preceding condition is equal to \(F\) (4.3.12).
\end{corollary}

\numpara{4.6.3}\label{IV.4.6.3}
We shall now concern ourselves with the case where \(X,R\)
are representable. Let us first introduce some terminology.
Besides conditions (a) to (d) introduced in 3.4.1, we shall use
other conditions on a family (M) of morphisms of \(\mathcal C\),
which we state below, recalling conditions (a) to (c) already
given, for completeness.

(a) (M) is stable under base extension.

(b) The composite of two elements of (M) is in (M).

(c) Every isomorphism is an element of (M).

\((\mathrm d_{\mathcal T})\) Every element of (M) is covering.
\nde{Recall that \(\mathcal T\) denotes the given topology on
\(\mathcal C\), less fine than the canonical topology.}

\((\mathrm e_{\mathcal T})\) Let \(f:X\to Y\) be a morphism of
\(\mathcal C\). If there exists a refinement \(R\) of \(Y\)
such that for every \(Y'\to R\), \(X\times_Y Y'\to Y'\)
is an element of (M), then \(f\) is an element of (M).

Recall that (a) and (b) imply

(a') The cartesian product of two elements of (M) is an element of (M).

On the other hand, (a) and \((\mathrm d_{\mathcal T})\)
imply, by 4.3.9:

(d') Every element of (M) is a universal effective epimorphism.

\numpara{4.6.4}\label{IV.4.6.4}
The preceding conditions are satisfied by the family of covering
morphisms, denoted \((\mathrm M_{\mathcal T})\), when
\(\mathcal C\) possesses fibered products.
Indeed (cf. 4.2.3), (a) follows from (C 1), (b) from (C 2),
(c) from (C 4), \((\mathrm d_{\mathcal T})\) from the definition,
and \((\mathrm e_{\mathcal T})\) from (C 5).
The results that we shall establish for a family satisfying
these conditions will apply in particular to the family
\((\mathrm M_{\mathcal T})\). In particular, one may take
for \(\mathcal T\) the canonical topology, and for (M)
the family of universal effective epimorphisms.

\begin{lemma}{4.6.5}\label{IV.4.6.5}
Let (M) be a family of morphisms satisfying the preceding
properties (a) to \((\mathrm e_{\mathcal T})\).
Let \(R\) be a \(\mathcal C\)-equivalence relation in
\(X\in\Ob\mathcal C\), of type (M).
Let \(\widetilde X\) be the sheaf defined by \(X\),
\(\widetilde R\) the \(\widetilde{\mathcal C}\)-equivalence
relation in \(\widetilde X\) defined by \(R\), and
\(\widetilde X/\widetilde R\) the quotient sheaf.
For \(R\) to be (M)-effective, it is necessary and sufficient
that \(\widetilde X/\widetilde R\) be representable.
If this is so, \(\widetilde X/\widetilde R\) is represented
by the quotient \(X/R\).
\end{lemma}
\begin{proof}
First suppose that \(R\) is (M)-effective and write \(Y=X/R\).
The canonical morphism \(p:X\to Y\) is an element of (M),
hence covering by \((\mathrm d_{\mathcal T})\).
The corresponding morphism
\[
\widetilde p:\widetilde X\longrightarrow
\widetilde X/\widetilde R
\]
is therefore a universal effective epimorphism of
\(\widetilde{\mathcal C}\) (4.4.3), hence identifies
\(\widetilde X/\widetilde R\) with the quotient of
\(\widetilde X\) by the equivalence relation \(R'\)
defined in \(\widetilde X\) by \(\widetilde p\).
Since the canonical functor \(\mathcal C\to
\widetilde{\mathcal C}\) commutes with fibered products,
\(R'\) is precisely \(\widetilde R\), since \(R\) is the
equivalence relation defined by \(p\).

Conversely, suppose that \(\widetilde X/\widetilde R\)
is represented by an object \(Y\) of \(\mathcal C\).
Let \(p:X\to Y\) be the morphism deduced from the canonical
morphism \(\widetilde X\to\widetilde X/\widetilde R\);
it is a covering morphism by 4.4.3.
It is clear, as before, that \(R\) is the equivalence relation
defined by \(p\). It remains only to show that \(p\in(\mathrm M)\).
Now the cartesian square
\[
\xymatrix{
R\ar[r]^{\sim}\ar[dr]_{p_2}
 & X\times_Y X\ar[r]\ar[d]_{\mathrm{pr}_2} & X\ar[d]^p\\
 & X\ar[r]_p & Y }
\]
shows that \(p\) becomes \(p_2\), which is an element of (M),
after base change by the covering morphism \(p\).
One concludes by \((\mathrm e_{\mathcal T})\).
\end{proof}

\begin{corollary}{4.6.5.1}\label{IV.4.6.5.1}
\nde{This corollary has been added.}
Let (M) satisfy the preceding properties (a) to
\((\mathrm e_{\mathcal T})\), and let \(f:G\to G'\) be a
morphism of \(\mathcal C\)-groups such that \(f\in(\mathrm M)\).
Suppose that \(\Ker(f)\) is representable (which is the
case if \(\mathcal C\) possesses a final object \(e\)).
Then the equivalence relation in \(G\) defined by \(H=\Ker(f)\)
is (M)-effective, and \(G'\) represents the quotient sheaf
\(\widetilde G/\widetilde H\) for the topology \(\mathcal T\).
\end{corollary}
\begin{proof}
This follows from 4.6.5 and 3.4.7.1.
\end{proof}
We are now in a position to state the principal result of this number.

\begin{theorem}{4.6.6}\label{IV.4.6.6}
Let (M) be a family of morphisms satisfying axioms (a) to
\((\mathrm e_{\mathcal T})\) of 4.6.3.
Let \(R\) be a \(\mathcal C\)-equivalence relation of type
(M) (cf. 3.4.3) in the object \(X\) of \(\mathcal C\).
Consider the functor \(F\in\Ob\widehat{\mathcal C}\)
defined as follows:
\[
F(S)=\left\{\begin{array}{l}
\text{sub-}S\text{-sheaves }Z\text{ of }X_S\text{ stable under }R\times S\\
\text{whose quotient by }R_Z\text{ is }S
\end{array}\right\}.
\]
Let \(F_0\) be the subfunctor of \(F\) defined as follows:
\(F_0(S)\) consists of the representable \(Z\in F(S)\), that is:
\[
F_0(S)=\left\{\begin{array}{l}
\text{sub-}\mathcal C_{/S}\text{-objects }Z\text{ of }X_S
\text{ stable under }R\times S,\\
\text{such that }R_Z\text{ is (M)-effective with quotient }S\\
\text{(i.e. such that }Z\to S\text{ belongs to (M) and }R_Z\simeq Z\times_S Z)
\end{array}\right\}.
\]
Then:

(i) The morphism \(p\in\Hom(X,F)=F(X)\) defined by the
subobject \(R\) of \(X\times X\) identifies \(F\) with the
quotient sheaf of \(X\) by \(R\).

(ii) The following conditions are equivalent:

a) \(F\) is representable.

b) \(F_0\) is representable.

c) \(R\) is (M)-effective.

Under these conditions, \(F=F_0=X/R\).

(iii) Let (N) be a family of morphisms stable under base
change, such that for every covering family \(\{S_i\to S\}\)
and every family \(\{T_i\to S_i\}\) of morphisms of (N),
every descent datum on the \(T_i\) relative to \(\{S_i\to S\}\)
is effective. Suppose that \(X\) is base-changeable (cf. 1.6.0)
and that the morphism \(R\to X\times X\) is an element of (N).
Then \(F_0=F\).
\end{theorem}
\begin{proof}
(i) has already been proved (4.6.1).

(ii) One has seen the equivalence of a) and c), as well as the
equality \(F=X/R\). It remains to prove that b) or c) implies
\(F_0=F\). First remark, as indeed asserted in the statement,
that \(F_0\) is a subfunctor of \(F\); indeed, for every
\(S\in\Ob\mathcal C\) and every \(Z\in F_0(S)\), the
morphism \(Z\to S\) is base-changeable, hence
\(Z\times_S S'\) is an element of \(F_0(S')\) for every
\(S'\to S\). Since \(R\in F(X)\) belongs to \(F_0(X)\),
4.6.2 shows that b) implies \(F_0=F\).

Now suppose that c) holds and let \(Q\) be an object of
\(\mathcal C\) representing \(X/R\).
Then the morphism \(X\to Q\) is an element of (M), and for
every \(S\in\Ob\mathcal C\) and every \(Z\in F(S)\),
diagram (*) of 4.6.1 shows that
\(Z=S\times_{(Q\times S)}(X\times S)\) is representable,
and \(Z\to S\) belongs to (M), hence \(Z\in F_0(S)\).

(iii) Let \(f\in\Hom(S,F)\) correspond to \(Z\in F(S)\).
One must show that \(f\) factors through \(F_0\), that is,
that \(Z\) is representable. This is first clear if \(f\)
factors through \(X\), by virtue of:

\begin{lemma}{4.6.7}\label{IV.4.6.7}
Let \(x_0\in X(S)\). The image of \(x_0\) in \(F(S)\)
corresponds to the subsheaf \(Z\) of \(X_S\) defined by
the two cartesian squares
\[
\xymatrix{
X_S\ar[r]^{\mathrm{id}_{X_S}\times x_0}
 & X_S\times_S X_S\ar[r] & X\times X\\
Z\ar[u]\ar[r] & R_S\ar[u]\ar[r] & R\ar[u] }.
\]
\end{lemma}
\begin{proof}
This lemma follows at once from the description of the
morphism \(X\to F\).
\end{proof}
Let us return to the proof of the theorem.
If \(f\) factors through \(X\), then \(Z\) is representable,
and, since \(R\to X\times X\) is an element of (N),
so is \(Z\to X_S\).
In general \(f\) does not necessarily factor through \(X\);
but, since \(X\to F\) is covering (4.4.3), there exists by
4.4.8 (vii) a covering family \(\{S_i\to S\}\), and for
each \(i\) a morphism \(S_i\to X\) making the diagram
\[
\xymatrix{X\ar[r] & F\\ S_i\ar[u]\ar[r] & S\ar[u]_f }
\]
commutative.
By the preceding argument, the morphism \(f_i:S_i\to F\)
defined by the preceding diagram belongs to \(\Hom(S_i,F_0)\)
and corresponds to the subsheaf \(Z\times_S S_i\) of \(X_{S_i}\).
The morphism \(Z\times_S S_i\to X_{S_i}\) is an element of (N),
and the family \(X_{S_i}\to X_S\) is covering.
It therefore remains only to establish:

\begin{proposition}{4.6.8}\label{IV.4.6.8}
Let \(\{S_i\to S\}\) be a covering family, and \(Z\) a sheaf
over \(S\). Suppose that for each \(i\), the \(S_i\)-functor
\(Z\times_S S_i\) is represented by an object \(T_i\).
Then the family of the \(T_i\) is endowed with a canonical
descent datum relative to \(S_i\to S\).
For \(Z\) to be representable, it is necessary and sufficient
that this datum be effective; if it is so, the
\og descended\fg{} object represents \(Z\).
\end{proposition}
\begin{proof}
First remark that by 4.4.3, \(S_i\to S\) is a universal
effective epimorphic family in \(\widetilde{\mathcal C}\),
hence a descent family in \(\widetilde{\mathcal C}\) (2.3).
If \(Z\) is represented by the object \(T\), then \(T\times_S S_i\)
(considered as a sheaf) is isomorphic to \(Z\times_S S_i\),
hence the descent datum on the \(T_i\) is effective, and
the descended object (unique) is isomorphic to \(Z\).
Conversely, suppose that the canonical descent datum on the
\(T_i\) is effective, and let \(T\) be the descended object.
Since the family \(\{S_i\to S\}\) is a descent family in
\(\widetilde{\mathcal C}\), there exists an \(S\)-morphism
\(T\to Z\) which by base extension to each \(S_i\) gives
back the canonical morphism \(T_i\to Z\times_S S_i\).
This morphism is locally an isomorphism; since \(T,Z\)
are sheaves, it follows from 4.5.8 that it is an isomorphism.
\end{proof}

\end{proof}
\begin{corollary}{4.6.9}\label{IV.4.6.9}
Let \(R\) be an (M)-effective equivalence relation in \(X\).
For every sheaf \(F\), the map
\[
\Hom(X/R,F)\longrightarrow\Hom(X,F)
\]
identifies the first set with the subset of the second consisting
of morphisms compatible with \(R\).
\end{corollary}
\begin{corollary}{4.6.10}\label{IV.4.6.10}
Let \(\mathcal T'\) be a topology less fine than \(\mathcal T\),
for which the morphisms of (M) are covering.
Under the conditions of 4.6.6 (iii), \(X/R\) is also the
quotient sheaf of \(X\) by \(R\) in every topology intermediate
between \(\mathcal T'\) and the canonical topology.
\end{corollary}
\begin{remark}{4.6.11}\label{IV.4.6.11}
If, in the statement of 4.6.6 (iii), one assumes moreover
that, under the hypotheses of the text, if one denotes by
\(T\) the descended object, the morphism \(T\to S\) is an
element of (N), then the inclusion morphisms \(Z\hookrightarrow X_S\)
are also elements of (N), as follows at once from the
construction of \(Z\) by descent.
\end{remark}
\begin{remark}{4.6.12}\label{IV.4.6.12}
The implications c) \(\Rightarrow\) b) \(\Rightarrow\) a)
and c) \(\Rightarrow[F_0=F=X/R]\) have been established
without recourse to the \og sufficient\fg{} part of Lemma 4.6.5,
which is the only place where condition \((\mathrm e_{\mathcal T})\)
is used. They therefore remain valid if (M) satisfies only
conditions (a) to \((\mathrm d_{\mathcal T})\).
An example of such a family (M) is that of base-changeable
covering morphisms (compare with 4.6.4).
In the case of the canonical topology, these are precisely
the universal effective epimorphisms. One therefore has:
\end{remark}

\begin{corollary}{4.6.13}\label{IV.4.6.13}
Let \(R\) be a universal effective equivalence relation in \(X\).
Then the object \(X/R\) of \(\mathcal C\) is the quotient sheaf
of \(X\) by \(R\) for the canonical topology.
It represents the following functor:
\((X/R)(S)\) is the set of sub-\(\mathcal C_{/S}\)-objects
\(Z\) of \(X_S\) stable under \(R\times S\) and such that
the induced equivalence relation is universal effective
with quotient \(S\).
\end{corollary}
Similarly, for an arbitrary topology:
\begin{corollary}{4.6.14}\label{IV.4.6.14}
Let (M) be the family of base-changeable covering morphisms.
If \(R\) is an (M)-effective equivalence relation in \(X\),
then the object \(X/R\) of \(\mathcal C\) is the quotient
sheaf of \(X\) by \(R\), and represents the functor \(F_0\)
of 4.6.6.
\end{corollary}

\begin{scholium}{4.6.15}\label{IV.4.6.15}
We may now add the following precisions to 4.5.5.
Whereas in questions involving exclusively projective limits
(fibered products, algebraic structures, etc.), one may,
by the results of Exposé I and 4.5.5, identify \(\mathcal C\)
indifferently with a full subcategory of
\(\widetilde{\mathcal C}\) or \(\widehat{\mathcal C}\),
the same is not true in questions mixing projective and
inductive limits. In all questions involving both projective
limits and inductive limits, in particular passages to
the quotient (example: group structure on the quotient of
a group by an invariant subgroup), we shall consider the
given category as embedded in the category of sheaves;
thus if \(R\) is a \(\mathcal C\)-equivalence relation in
the object \(X\) of \(\mathcal C\), \(X/R\) will denote
the quotient sheaf of \(X\) by \(R\) (previously denoted
by \(j(X)/j(R)\)), hence, in the case where this sheaf is
representable, the object which it represents.
The preceding results show that in the most important cases
a quotient in \(\mathcal C\) will also be a quotient in the
category of sheaves; in any event we forbid ourselves the
use of the notation \(X/R\) for a quotient in \(\mathcal C\)
which does not coincide with the quotient in
\(\widetilde{\mathcal C}\) (for example, one which is not
universal), thus modifying the definitions of no.\ 3.

To study a problem of the above type, one therefore first
works in the category of sheaves, where all the customary
results are valid (cf. no.\ 4.4), then specializes the
results obtained to the initial category, using the results
of this number and, when available, criteria for effectiveness
of descent. We shall see examples of this method in the
following numbers.
\end{scholium}

\sgasection{4.7}{Use of effectiveness criteria: isomorphism theorem}
In this number we give an example of the use of effectiveness
criteria. The initial data are a topology \(\mathcal T\)
on \(\mathcal C\) (still less fine than the canonical topology),
a family (M) of morphisms of \(\mathcal C\) satisfying axioms
(a) to \((\mathrm e_{\mathcal T})\) of 4.6.3, and a family
(N) of morphisms of \(\mathcal C\) which may satisfy the following axioms:

(a) (N) is stable under base extension.

\((\mathrm f_{\mathcal T})\) \og Morphisms of (N) descend
under the given topology\fg; that is: for every
\(S\in\Ob\mathcal C\), every covering family \(\{S_i\to S\}\)
and every family \(\{T_i\to S_i\}\) of morphisms of (N),
every descent datum on the \(T_i\) relative to \(\{S_i\to S\}\)
is effective, and if one denotes by \(T\) the descended object,
the morphism \(T\to S\) is an element of (N).

Since every element of (M) is covering (condition
4.6.3 \((\mathrm d_{\mathcal T})\)),
\((\mathrm f_{\mathcal T})\) implies the following axiom:
\nde{Axiom \((\mathrm f_{\mathrm M})\) (which appeared before
Proposition 4.7.2) has been placed here.}

\((\mathrm f_{\mathrm M})\) If \(Y\to X\) is an element
of (N) and \(X\to X_1\) an element of (M), every descent
datum on \(Y\) relative to \(X\to X_1\) is effective;
if one denotes by \(Y_1\) the descended object, \(Y_1\to X_1\)
is an element of (N).

Let us immediately point out an example of this situation,
which will be treated later: \(\mathcal C\) is the category
of schemes, \(\mathcal T\) the faithfully flat quasi-compact
topology; (M) the family of faithfully flat quasi-compact
morphisms, (N) the family of closed immersions or that
of quasi-compact immersions.
\nde{Cf. \(\S\) 6.4, see also VIA, 5.3.1.}

Recall the principal result of 4.6.6 (taking account of 4.6.11):
\begin{proposition}{4.7.1}\label{IV.4.7.1}
If \(X\) is a base-changeable object of \(\mathcal C\),
\(R\) an equivalence relation of type (M) in \(X\)
such that \(R\to X\times X\) is an element of (N),
with (N) satisfying (a) and \((\mathrm f_{\mathcal T})\),
then the quotient sheaf \(X/R\) is defined by
\[
(X/R)(S)=\left\{\begin{array}{l}
\text{sub-}S\text{-objects }Z\text{ of }X_S\text{ stable under }R\times S,\\
\text{such that }Z\to X_S\text{ belongs to (N), }Z\to S\text{ is covering}\\
\text{(or an element of (M)), and }R_Z\simeq Z\times_S Z
\end{array}\right\}.
\]
\end{proposition}
Moreover, one has:
\begin{proposition}{4.7.2}\label{IV.4.7.2}
Let \(X\in\Ob\mathcal C\), and \(R\) an (M)-effective
equivalence relation in \(X\). Let (N) be a family of
morphisms satisfying (a) and \((\mathrm f_{\mathrm M})\).

For every subobject \(Y\) of \(X\) stable under \(R\)
and such that \(Y\to X\) belongs to (N), the equivalence
relation induced in \(Y\) by \(R\) is (M)-effective, and
the quotient \(Y/R_Y=Y'\) is a subobject of \(X'=X/R\)
such that \(Y'\to X'\) belongs to (N).

The map \(Y\mapsto Y'\) is a bijection between the set
of subobjects \(Y\) of \(X\) stable under \(R\), such
that \(Y\to X\) belongs to (N), and the set of subobjects
\(Y'\) of \(X'\) such that \(Y'\to X'\) belongs to (N).
The inverse map is \(Y'\mapsto Y'\times_{X'}X\).
\end{proposition}
\begin{proof}
Since \(R\) is (M)-effective, the morphism \(X\to X'\)
belongs to (M). Let \(Y'\) be a subobject of \(X'\)
such that the canonical morphism \(Y'\to X'\) belongs
to (N). Then the subobject \(Y=Y'\times_{X'}X\) of
\(X\) is stable under \(R\), and the morphism \(Y\to X\)
(resp. \(Y\to Y'\)) belongs to (N) (resp. (M)), since
(N) and (M) are stable under base change.
Denote by \(R_Y\) the equivalence relation induced in
\(Y\) by \(R\). By 4.4.11, the quotient sheaf \(Y/R_Y\)
is represented by \(Y'\), and hence, by 4.6.5,
\(R_Y\) is (M)-effective.

Conversely, let us show that every subobject \(Y\) of
\(X\) stable under \(R\), such that the structural
morphism \(Y\to X\) belongs to (N), is obtained in this way.
Indeed, if \(Y\) is stable under \(R\), its two images
in \(R=X\times_{X'}X\) are identical, and \(Y\) is
endowed with a descent datum relative to \(X\to X'\);
the desired result follows, since the family (N)
satisfies axiom \((\mathrm f_{\mathrm M})\).
\end{proof}
